%%%PAGE 282 rot=0 legibility=fair summary=贴有两道印刷题并带手写解答：弹射器与竖直圆轨道（求弹性势能范围，答0.2J≤Ep<0.23J或0.33J<Ep<0.43J）；小车与墙壁多次弹性碰撞（碰撞次数、v0范围30~34m/s、v0=28m/s时总冲量88/3N·s）
%%%BLOCK id=282-01 ch=06 sec=功能关系·弹射器与竖直圆轨道 kind=problem alt=04 cont=none y=0.00-0.49
% 贴纸题（印刷体）+手写解答。照片上边缘切掉了题干的前一两行（题号与开头无法辨认）；原稿：m=0.01kg、R=0.8m、r=0.1m、μ=0.5、L=2m、E_pm=0.5J均被框出，“滑块”“其余各部分轨道均光滑”被圈，“滑块与弹性板作用后以等大速率弹回”处有一道横线穿过（疑为划线强调），“最终静止在轨道FG中点的左侧区域内”下有划线
\begin{problem}
\unclear{…}安装在水平\unclear{…}% 第一行被照片上边缘切断，仅见右端“安装在水平”几字

\unclear{…}射器、水平直轨道$AB$、圆心为$O_1$的竖直半圆轨道$BCD$、圆心为$O_2$的竖直半圆管道$DEF$、水平直轨道$FG$及弹性板等组成，轨道各部分平滑连接。已知滑块（可视为质点）质量$m=0.01\unit{kg}$，轨道$BCD$的半径$R=0.8\unit{m}$，管道$DEF$的半径$r=0.1\unit{m}$，滑块与轨道$FG$间的动摩擦因数$\mu=0.5$，其余各部分轨道均光滑，轨道$FG$的长度$L=2\unit{m}$，弹射器中弹簧的弹性势能最大值$E_{pm}=0.5\unit{J}$，滑块与弹簧作用后，弹簧的弹性势能完全转化为滑块的动能，滑块与弹性板作用后以等大速率弹回，$g=10\unit{m/s^2}$。

%FIG p282_a 0.68 0.03 0.99 0.185
\notefig[0.3]{figures/p282_a.png}
（1）若滑块被弹簧弹开后恰能沿轨道$BCD$到达$D$点，求弹簧的弹性势能$E_p$；

（2）若滑块在运动过程中最终静止在轨道$FG$中点的左侧区域内，求弹簧的弹性势能$E_p$的范围
\begin{solution}
（1）$mg=\dfrac{mv_{\min}^2}{R}$

$E_p=mg\,2R+\dfrac{1}{2}mv_{\min}^2$

$E_p=0.2\unit{J}$

（2）若滑块在运动过程中第一次恰运动到$F$点

$E_{p\min}=mg(2R+2r)=0.18\unit{J}<0.2\unit{J}$\qquad $E_{p\min}=0.2\unit{J}$

运动过程中\ 若恰好第一次运动到$FG$中点。\ $E_{p\max}=mg(2R+2r)+\mu mg\dfrac{L}{2}=0.23\unit{J}$% “运动过程中”为行首上方的小字（插入）

第二次\quad $E_{p\min}=mg(2R+2r)+\mu mg\dfrac{3L}{2}=0.33\unit{J}$

第三次\quad $E_{p\max}=mg(2R+2r)+\mu mg\dfrac{5}{2}L=0.43\unit{J}$

第四次\quad $E_{p\max}=mg(2R+2r)+\mu mg\dfrac{7}{2}L=0.53\unit{J}$

由$E_{p\max}=0.5\unit{J}$

综上 弹性势能范围：$0.2\unit{J}\le E_p<0.23\unit{J}$ 或 $0.33\unit{J}<E_p<0.43\unit{J}$
\end{solution}
\end{problem}
%%%END
%%%BLOCK id=282-02 ch=07 sec=动量·小车与墙壁多次弹性碰撞 kind=problem alt=03,06 cont=none y=0.46-1.00
% 贴纸题（印刷体）+手写解答（含红笔）。原稿：x0=2m、m1=1kg、μ=0.2、m2=2kg均被框出，“光滑”“v0”被圈，“弹性碰撞”下有横线，“且能发生多次碰撞”被圈；(1)(2)问末尾有红笔小字“3”“5”（疑为分值）；(3)中“v0=28m/s”上方红笔写“28m/s”，问句下有红色双下划线
\begin{problem}[21，15分]
如图所示，足够长的小车静放在光滑水平面上，车右端与墙壁相距为$x_0=2\unit{m}$，在小车左端放一个质量为$m_1=1\unit{kg}$的小物体（可视为质点），与小车的动摩擦因数为$\mu=0.2$，小车的质量为$m_2=2\unit{kg}$，现给小物体一个水平向右的初始速度$v_0$，已知小车与墙壁的碰撞始终为弹性碰撞，且能发生多次碰撞，$g=10\unit{m/s^2}$。

%FIG p282_b 0.64 0.497 0.93 0.581
\notefig[0.3]{figures/p282_b.png}
（1）求小车第一次碰墙前的速度和该次碰撞墙壁给小车的冲量大小\ \red{3}

（2）在整个的运动过程中，若墙壁给了小车$I=32\unit{N\cdot s}$的冲量，求小车与墙壁碰撞的次数和$v_0$的范围\ \red{5}

（3）如果$v_0=28\unit{m/s}$，墙壁给了小车多少冲量？最终物体在小车上滑行的路程多大？\ \red{28$\unit{m/s}$}% 红笔“28m/s”写在“v0=28m/s”上方
\begin{solution}
（1）由题意运动过程不能共速

$a=\mu g=1\unit{m/s^2}$\quad $v_2=\sqrt{2a_2x_0}$% CHECK: μg=2m/s²；原稿写 a=μg=1m/s²（应指小车加速度 a2=μm1g/m2=1m/s²），按原样转录

（2）$n=\dfrac{I}{I_2}=4$次

为保证小车与墙壁碰撞4次，分为两种临界情况。

①当$v_0$取最小值时，滑块与小车恰在第四次碰撞共速

$t_2=\dfrac{7v_2}{a_2}=14\unit{s}$

对滑块\ $a_1=\mu g=2\unit{m/s^2}$\quad \red{$v_0-28=2$\quad $v_0=30$}\quad \red{$v_0-32=2$\quad $v_0=34$}

$v_2=v_0-a_1t_2$\quad $\therefore v_0=30\unit{m/s}$

②当$v_0$取最大值，滑块与小车在$v=0$时共速\quad $t_2'=\dfrac{8v_2}{a_2}=16\unit{s}$

对滑块\ $v_2=v_0-a_1t_2$，$v_0'=34\unit{m/s}$\quad 故$30\unit{m/s}\le v_0\le34\unit{m/s}$

（3）

%FIG p282_c 0.44 0.676 0.84 0.80
\notefig[0.5]{figures/p282_c.png}
\red{$v_{\text{块}}=28-2t$}

$a_{\text{板}}=1\unit{m/s^2}$\quad $a_{\text{块}}=2\unit{m/s^2}$

碰4次之后\ 小车$v_2''=\dfrac{4}{3}\unit{m/s}$\ 向左

滑块$v_1''=\dfrac{4}{3}\unit{m/s}$\ 向右

$m_1v_1''-m_2v_2''=(m_1+m_2)v_3'$

$v_3'=-\dfrac{4}{9}\unit{m/s}$

$\mu m_1gs=\dfrac{1}{2}m_1v_0^2-\dfrac{1}{2}(m_1+m_2)v_3'^2$

$s=\dfrac{5288}{27}\unit{m}$

\red{$\leftarrow$ 这是1次碰撞}% 红笔小箭头，左侧另有一道红色斜划线

碰3次 第4次碰前共速，且速度向右

前3次\quad $I_2'=3I_2=24\unit{N\cdot s}$

第3次碰后\quad 小车$v_2'=2\unit{m/s}$\quad 滑块$v_1'=8\unit{m/s}$

$m_1v_1'-m_2v_2'=(m_1+m_2)v_3$

$v_3=\dfrac{4}{3}\unit{m/s}$\ 向右

小车与墙碰后\quad $I_3=2m_2v_3=\dfrac{16}{3}\unit{N\cdot s}$

全程墙对车冲量\quad $I=\dfrac{88}{3}\unit{N\cdot s}$
\end{solution}
\end{problem}
%%%END
%%%PAGEEND 282
